\documentclass[10pt,letterpaper]{article}
\usepackage[T1]{fontenc}
\usepackage[utf8]{inputenc}
\usepackage[spanish,es-tabla]{babel}
\usepackage{csquotes}
\usepackage{mathpazo}
\usepackage[scaled=0.9]{helvet}
\usepackage{courier}
\usepackage{calc}
\usepackage[top=2cm, bottom=2cm, left=2cm, right=2cm]{geometry}
\usepackage{xcolor}
\usepackage{booktabs}
\usepackage{longtable}
\usepackage{array}
\usepackage{ragged2e}
\usepackage{xurl}
\usepackage[strings]{underscore}
\usepackage{fancyhdr}
\usepackage{sectsty}
\usepackage{tocloft}
\usepackage[colorlinks=true,linkcolor=blue,urlcolor=blue,citecolor=blue,breaklinks=true]{hyperref}
\usepackage{enumitem}

\definecolor{uncpblue}{HTML}{1A237E}
\allsectionsfont{\color{uncpblue}}

\renewcommand{\cftsecleader}{\cftdotfill{\cftsecdotsep}}
\cftpagenumbersoff{section}
\cftpagenumbersoff{subsection}
\cftpagenumbersoff{subsubsection}
\renewcommand{\cftsecfont}{\normalfont}
\renewcommand{\cftsubsecfont}{\normalfont}
\renewcommand{\cftsubsubsecfont}{\normalfont}
\setlength{\parindent}{0pt}
\setlength{\parskip}{4pt}
\setlength{\tabcolsep}{3pt}
\renewcommand{\arraystretch}{0.85}

\pagestyle{fancy}
\fancyhf{}
\fancyhead[L]{\small\textcolor{uncpblue}{\textbf{SGSI-UNCP}}}
\fancyhead[R]{\small\textcolor{uncpblue}{D-SGSI-12}}
\fancyfoot[C]{\thepage}
\renewcommand{\headrulewidth}{0.4pt}
\renewcommand{\footrulewidth}{0.4pt}
\setlength{\headheight}{14pt}

\setlist{nosep,leftmargin=*,after=\vspace{2pt}}
\setlength{\emergencystretch}{2em}

\newcommand{\makecover}{
\begin{titlepage}
\centering
\vspace*{2cm}
{\Huge\bfseries\color{uncpblue} SGSI-UNCP\par}
\vspace{0.7cm}
{\LARGE Estrategia de Certificación ISO 27001:2022\par}
\vspace{0.5cm}
{\large Documento Base del Sistema de Gestión\\de Seguridad de la Información\par}
\vspace{1.5cm}
{\small Universidad Nacional del Centro del Perú\par}
\vfill
\end{titlepage}
}

\begin{document}

\makecover
\setcounter{tocdepth}{2}
\RaggedRight
\tableofcontents
\newpage

\RaggedRight
\section{Estrategia de Certificación ISO
27001:2022}\label{estrategia-de-certificaciuxf3n-iso-270012022}

\textbf{Organización:} Universidad Nacional del Centro del Perú (UNCP)
\textbf{Periodo:} 2026-2027 \textbf{Responsable:} Oficial de Seguridad y
Confianza Digital \textbf{Presupuesto Estimado:} S/ 120,000 -- 150,000
(auditoría de certificación + pre-auditoría + consultoría de apoyo)

\begin{center}\rule{0.5\linewidth}{0.5pt}\end{center}

\subsection{Objetivo}\label{objetivo}

Establecer la hoja de ruta, los criterios de selección de entidad
certificadora y las actividades preparatorias necesarias para que la
UNCP obtenga la certificación ISO/IEC 27001:2022, demostrando la
conformidad de su Sistema de Gestión de Seguridad de la Información.

\subsection{Selección de Entidad
Certificadora}\label{selecciuxf3n-de-entidad-certificadora}

\subsubsection{Criterios de Selección}\label{criterios-de-selecciuxf3n}

{\setlength{\RaggedRightRightskip}{0pt plus 5em}\small\sloppy \begin{longtable}[]{@{}
  >{\hspace{0pt}\RaggedRight\arraybackslash}p{(\linewidth - 4\tabcolsep) * \real{0.3333}}
  >{\raggedleft\arraybackslash}p{(\linewidth - 4\tabcolsep) * \real{0.3333}}
  >{\hspace{0pt}\RaggedRight\arraybackslash}p{(\linewidth - 4\tabcolsep) * \real{0.3333}}@{}}
\toprule\noalign{}
\begin{minipage}[b]{\linewidth}\raggedright
Criterio
\end{minipage} & \begin{minipage}[b]{\linewidth}\raggedleft
Peso
\end{minipage} & \begin{minipage}[b]{\linewidth}\raggedright
Descripción
\end{minipage} \\
\midrule\noalign{}
\endhead
\bottomrule\noalign{}
\endlastfoot
\textbf{Acreditación internacional} & 25\% & La entidad debe estar
acreditada por IAF/DAkkS/ANAB/UKAS. Sin acreditación, el certificado no
tiene validez internacional. \\
\textbf{Presencia en Perú} & 20\% & Debe contar con auditores locales o
capacidad de despliegue sin costos excesivos de viaje. \\
\textbf{Experiencia en el sector público/educativo} & 20\% & Experiencia
comprobable en certificación de universidades públicas o instituciones
del Estado peruano. \\
\textbf{Reputación y reconocimiento} & 15\% & Reconocimiento en el
mercado peruano y latinoamericano. \\
\textbf{Costo total (Etapa 1 + Etapa 2 + 3 años de vigilancia)} & 20\% &
Relación costo-beneficio del servicio completo. \\
\end{longtable}\normalsize}

\subsubsection{Entidades Certificadoras
Recomendadas}\label{entidades-certificadoras-recomendadas}

{\setlength{\RaggedRightRightskip}{0pt plus 5em}\small\sloppy \begin{longtable}[]{@{}
  >{\hspace{0pt}\RaggedRight\arraybackslash}p{(\linewidth - 10\tabcolsep) * \real{0.1739}}
  >{\hspace{0pt}\RaggedRight\arraybackslash}p{(\linewidth - 10\tabcolsep) * \real{0.1304}}
  >{\hspace{0pt}\RaggedRight\arraybackslash}p{(\linewidth - 10\tabcolsep) * \real{0.1739}}
  >{\hspace{0pt}\RaggedRight\arraybackslash}p{(\linewidth - 10\tabcolsep) * \real{0.1739}}
  >{\raggedleft\arraybackslash}p{(\linewidth - 10\tabcolsep) * \real{0.1739}}
  >{\hspace{0pt}\RaggedRight\arraybackslash}p{(\linewidth - 10\tabcolsep) * \real{0.1739}}@{}}
\toprule\noalign{}
\begin{minipage}[b]{\linewidth}\raggedright
Entidad
\end{minipage} & \begin{minipage}[b]{\linewidth}\raggedright
Acreditación
\end{minipage} & \begin{minipage}[b]{\linewidth}\raggedright
Presencia en Perú
\end{minipage} & \begin{minipage}[b]{\linewidth}\raggedright
Experiencia Sector Público
\end{minipage} & \begin{minipage}[b]{\linewidth}\raggedleft
Estimado (S/)
\end{minipage} & \begin{minipage}[b]{\linewidth}\raggedright
Ventaja Clave
\end{minipage} \\
\midrule\noalign{}
\endhead
\bottomrule\noalign{}
\endlastfoot
\textbf{AENOR} & ENAC (IAF) & Sí, oficina en Lima & Alta (múltiples
municipios y gobiernos regionales) & 120,000 - 140,000 & Amplia
experiencia en el sector público peruano \\
\textbf{BSI} & UKAS (IAF) & Sí, oficina en Lima & Alta (reconocimiento
global, creadores del estándar) & 130,000 - 150,000 & Prestigio
internacional; creadores de ISO 27001 \\
\textbf{SGS} & ANAB (IAF) & Sí, oficina en Lima & Media (más presencia
en sector privado) & 110,000 - 130,000 & Red global de auditores; buena
relación costo-beneficio \\
\textbf{ICONTEC} & ONAC (IAF) & Sí, representación en Perú & Alta (muy
fuerte en educación superior en Latinoamérica) & 100,000 - 120,000 &
Mayor experiencia en universidades latinoamericanas \\
\end{longtable}\normalsize}

\subsubsection{Proceso de Contratación
Recomendado}\label{proceso-de-contrataciuxf3n-recomendado}

\begin{enumerate}
\def\labelenumi{\arabic{enumi}.}
\tightlist
\item
  Enviar \textbf{solicitud de cotización (RFP)} a las 4 entidades en
  Q3-2026.
\item
  Evaluar las propuestas según los criterios de la sección anterior.
\item
  Solicitar referencias de al menos 2 instituciones similares
  certificadas por la entidad.
\item
  Entrevistar al equipo auditor propuesto para la Etapa 1.
\item
  Seleccionar y contratar en Q4-2026.
\end{enumerate}

\begin{center}\rule{0.5\linewidth}{0.5pt}\end{center}

\subsection{Proceso de Certificación
(Fases)}\label{proceso-de-certificaciuxf3n-fases}

\subsubsection{Cronograma General}\label{cronograma-general}

{\setlength{\RaggedRightRightskip}{0pt plus 5em}\small\sloppy \begin{longtable}[]{@{}
  >{\hspace{0pt}\RaggedRight\arraybackslash}p{(\linewidth - 8\tabcolsep) * \real{0.1905}}
  >{\hspace{0pt}\RaggedRight\arraybackslash}p{(\linewidth - 8\tabcolsep) * \real{0.1905}}
  >{\raggedleft\arraybackslash}p{(\linewidth - 8\tabcolsep) * \real{0.1905}}
  >{\centering\arraybackslash}p{(\linewidth - 8\tabcolsep) * \real{0.2381}}
  >{\hspace{0pt}\RaggedRight\arraybackslash}p{(\linewidth - 8\tabcolsep) * \real{0.1905}}@{}}
\toprule\noalign{}
\begin{minipage}[b]{\linewidth}\raggedright
Fase
\end{minipage} & \begin{minipage}[b]{\linewidth}\raggedright
Actividad
\end{minipage} & \begin{minipage}[b]{\linewidth}\raggedleft
Fecha Sugerida
\end{minipage} & \begin{minipage}[b]{\linewidth}\centering
Duración Estimada
\end{minipage} & \begin{minipage}[b]{\linewidth}\raggedright
Responsable
\end{minipage} \\
\midrule\noalign{}
\endhead
\bottomrule\noalign{}
\endlastfoot
\textbf{Preparación} & Auditoría de brecha (gap analysis) con consultora
externa & Q4 2026 & 2 semanas & Consultora externa + Oficial de
Seguridad \\
\textbf{Preparación} & Cierre de brechas identificadas & Q1 2027 & 3
meses & OTI + Oficial de Seguridad \\
\textbf{Preparación} & Micro-auditorías internas de verificación & Q1
2027 & Continuo & Equipo auditor interno \\
\textbf{Etapa 1} & Auditoría documental (revisión de diseño del SGSI) &
Q1 2027 (marzo) & 2 días & Entidad certificadora \\
\textbf{Preparación} & Cierre de no conformidades de Etapa 1 (si las
hubiera) & Q2 2027 & 2 meses & OTI + Oficial de Seguridad \\
\textbf{Etapa 2} & Auditoría de campo (implementación y evidencias) & Q3
2027 (agosto) & 4 días (2 auditores) & Entidad certificadora \\
\textbf{Cierre} & Cierre de no conformidades de Etapa 2 & Q3 2027 & 1
mes & OTI + Oficial de Seguridad \\
\textbf{Certificación} & Emisión del certificado ISO 27001:2022 & Q3/Q4
2027 & -- & Entidad certificadora \\
\textbf{Vigilancia} & Auditorías de vigilancia anuales & 2028, 2029,
2030 & 2 días cada una & Entidad certificadora \\
\end{longtable}\normalsize}

\subsubsection{Auditoría de Etapa 1 (Revisión
Documental)}\label{auditoruxeda-de-etapa-1-revisiuxf3n-documental}

\begin{itemize}
\tightlist
\item
  \textbf{Objetivo:} Verificar que el SGSI está correctamente diseñado y
  documentado conforme a ISO 27001:2022.
\item
  \textbf{Alcance:} Revisión de todos los documentos del SGSI: contexto,
  política, alcance, SoA, metodología de riesgos, procedimientos
  operativos.
\item
  \textbf{Duración típica:} 1-2 días con 1 auditor.
\item
  \textbf{Posibles hallazgos:} Documentación incompleta, falta de
  evidencias de implementación inicial, inconsistencias en la SoA.
\item
  \textbf{Si hay NC mayores:} No se puede pasar a Etapa 2 hasta
  cerrarlas.
\end{itemize}

\subsubsection{Auditoría de Etapa 2 (Verificación en
Campo)}\label{auditoruxeda-de-etapa-2-verificaciuxf3n-en-campo}

\begin{itemize}
\tightlist
\item
  \textbf{Objetivo:} Verificar que el SGSI está implementado, operativo
  y es eficaz.
\item
  \textbf{Alcance:} Entrevistas con personal clave, revisión de
  registros y evidencias (logs, actas, informes), verificación de
  controles técnicos y físicos.
\item
  \textbf{Duración típica:} 3-5 días con 2 auditores (dependiendo del
  tamaño y complejidad de la UNCP).
\item
  \textbf{Áreas a auditar (muestra):}

  \begin{itemize}
  \tightlist
  \item
    Rectorado/Comité de Gobierno Digital (liderazgo, revisión por la
    dirección).
  \item
    Oficina del Oficial de Seguridad (gestión de riesgos, SoA,
    métricas).
  \item
    OTI (gestión de accesos, operaciones, respuesta a incidentes,
    respaldos).
  \item
    Una o dos facultades (aplicación de controles en el día a día).
  \item
    Datacenter (seguridad física y ambiental).
  \end{itemize}
\end{itemize}

\subsubsection{Emisión del
Certificado}\label{emisiuxf3n-del-certificado}

\begin{itemize}
\tightlist
\item
  El certificado se emite tras el cierre de todas las no conformidades
  mayores y la aceptación del plan de corrección de las menores.
\item
  \textbf{Vigencia:} 3 años, con auditorías de vigilancia anuales.
\item
  \textbf{Alcance del certificado:} Debe coincidir exactamente con el
  Alcance del SGSI definido en D-SGSI-02.
\end{itemize}

\subsubsection{Auditorías de Vigilancia (Años 2, 3 y
4)}\label{auditoruxedas-de-vigilancia-auxf1os-2-3-y-4}

\begin{itemize}
\tightlist
\item
  \textbf{Año 2 (2028):} Auditoría de vigilancia 1 (énfasis en mejora
  continua y cierre de hallazgos previos).
\item
  \textbf{Año 3 (2029):} Auditoría de vigilancia 2 (muestra de controles
  diferente al año anterior).
\item
  \textbf{Año 4 (2030):} Auditoría de recertificación (equivalente a
  Etapa 1 + Etapa 2).
\end{itemize}

\begin{center}\rule{0.5\linewidth}{0.5pt}\end{center}

\subsection{Preparación Final
(Pre-auditoría)}\label{preparaciuxf3n-final-pre-auditoruxeda}

\subsubsection{Auditoría de Brecha (Gap
Analysis)}\label{auditoruxeda-de-brecha-gap-analysis}

Se recomienda contratar a una consultora externa (no la misma entidad
certificadora) para realizar una \textbf{auditoría de brecha} en Q4-2026
que evalúe:

{\setlength{\RaggedRightRightskip}{0pt plus 5em}\small\sloppy \begin{longtable}[]{@{}
  >{\hspace{0pt}\RaggedRight\arraybackslash}p{(\linewidth - 2\tabcolsep) * \real{0.5000}}
  >{\hspace{0pt}\RaggedRight\arraybackslash}p{(\linewidth - 2\tabcolsep) * \real{0.5000}}@{}}
\toprule\noalign{}
\begin{minipage}[b]{\linewidth}\raggedright
Aspecto
\end{minipage} & \begin{minipage}[b]{\linewidth}\raggedright
Qué se Revisa
\end{minipage} \\
\midrule\noalign{}
\endhead
\bottomrule\noalign{}
\endlastfoot
\textbf{Documentación} & ¿Todos los documentos requeridos por ISO 27001
existen y están completos? \\
\textbf{Implementación} & ¿Los controles declarados en la SoA están
realmente implementados? \\
\textbf{Evidencias} & ¿Existen registros que demuestren la operación del
SGSI? \\
\textbf{Personal} & ¿El personal conoce sus roles y responsabilidades en
el SGSI? \\
\textbf{Riesgos} & ¿La evaluación de riesgos es completa y está
actualizada? \\
\end{longtable}\normalsize}

\subsubsection{Cierre de Brechas}\label{cierre-de-brechas}

Los hallazgos de la auditoría de brecha se registran como RACs
(\textbf{F-SGSI-04}) y se les asigna responsable y fecha límite. El
Oficial de Seguridad realiza el seguimiento semanal hasta el cierre
total.

\begin{center}\rule{0.5\linewidth}{0.5pt}\end{center}

\subsection{Recursos Necesarios}\label{recursos-necesarios}

{\setlength{\RaggedRightRightskip}{0pt plus 5em}\small\sloppy \begin{longtable}[]{@{}
  >{\hspace{0pt}\RaggedRight\arraybackslash}p{(\linewidth - 2\tabcolsep) * \real{0.5000}}
  >{\raggedleft\arraybackslash}p{(\linewidth - 2\tabcolsep) * \real{0.5000}}@{}}
\toprule\noalign{}
\begin{minipage}[b]{\linewidth}\raggedright
Concepto
\end{minipage} & \begin{minipage}[b]{\linewidth}\raggedleft
Estimado (S/)
\end{minipage} \\
\midrule\noalign{}
\endhead
\bottomrule\noalign{}
\endlastfoot
Auditoría de brecha (consultora externa) & 25,000 -- 35,000 \\
Honorarios de certificación (Etapa 1 + Etapa 2 + 3 años de vigilancia) &
100,000 -- 120,000 \\
Adecuaciones técnicas menores (derivadas de la auditoría de brecha) &
20,000 -- 50,000 \\
Capacitación adicional del personal (si se detectan brechas de
competencia) & 10,000 -- 20,000 \\
\textbf{Total} & \textbf{155,000 -- 225,000} \\
\end{longtable}\normalsize}

\begin{center}\rule{0.5\linewidth}{0.5pt}\end{center}

\subsection{Riesgos Identificados para la
Certificación}\label{riesgos-identificados-para-la-certificaciuxf3n}

{\setlength{\RaggedRightRightskip}{0pt plus 5em}\small\sloppy \begin{longtable}[]{@{}
  >{\hspace{0pt}\RaggedRight\arraybackslash}p{(\linewidth - 6\tabcolsep) * \real{0.2222}}
  >{\centering\arraybackslash}p{(\linewidth - 6\tabcolsep) * \real{0.2778}}
  >{\centering\arraybackslash}p{(\linewidth - 6\tabcolsep) * \real{0.2778}}
  >{\hspace{0pt}\RaggedRight\arraybackslash}p{(\linewidth - 6\tabcolsep) * \real{0.2222}}@{}}
\toprule\noalign{}
\begin{minipage}[b]{\linewidth}\raggedright
Riesgo
\end{minipage} & \begin{minipage}[b]{\linewidth}\centering
Probabilidad
\end{minipage} & \begin{minipage}[b]{\linewidth}\centering
Impacto
\end{minipage} & \begin{minipage}[b]{\linewidth}\raggedright
Mitigación
\end{minipage} \\
\midrule\noalign{}
\endhead
\bottomrule\noalign{}
\endlastfoot
Documentación incompleta en Etapa 1 & Baja & Alto & Realizar auditoría
de brecha previa; checklist de documentos requeridos \\
Personal no preparado para entrevistas de Etapa 2 & Media & Alto &
Realizar simulacros de auditoría con el personal clave \\
Hallazgos técnicos no corregidos a tiempo & Media & Medio & Establecer
un plan de acción con hitos quincenales \\
Cambios en la normativa peruana que afecten los requisitos & Baja &
Medio & Monitoreo mensual del marco legal (D-SGSI-00) \\
\end{longtable}\normalsize}

\begin{center}\rule{0.5\linewidth}{0.5pt}\end{center}

\subsection{Documentos Relacionados}\label{documentos-relacionados}

{\setlength{\RaggedRightRightskip}{0pt plus 5em}\small\sloppy \begin{longtable}[]{@{}
  >{\hspace{0pt}\RaggedRight\arraybackslash}p{(\linewidth - 2\tabcolsep) * \real{0.5000}}
  >{\hspace{0pt}\RaggedRight\arraybackslash}p{(\linewidth - 2\tabcolsep) * \real{0.5000}}@{}}
\toprule\noalign{}
\begin{minipage}[b]{\linewidth}\raggedright
Código
\end{minipage} & \begin{minipage}[b]{\linewidth}\raggedright
Nombre
\end{minipage} \\
\midrule\noalign{}
\endhead
\bottomrule\noalign{}
\endlastfoot
D-SGSI-02 & Alcance del SGSI \\
D-SGSI-05 & Declaración de Aplicabilidad (SoA) \\
P-SGSI-09 & Metodología de Auditoría Basada en Riesgos \\
R-SGSI-04 & Programa Maestro de Auditoría \\
INFORME-FINAL-SGSI-UNCP & Informe Final Ejecutivo del Diseño del SGSI \\
\end{longtable}\normalsize}

\begin{center}\rule{0.5\linewidth}{0.5pt}\end{center}

\textbf{Elaborado por:} Oficial de Seguridad y Confianza Digital -- UNCP

\end{document}
